\section{Related Work}
\subsection{Learning in radar--jammer games}
Recent radar anti-jamming studies have used multi-agent reinforcement learning and game-theoretic formulations to adapt waveform, frequency, and power decisions. Neural fictitious self-play (NFSP) has been applied to imperfect-information games, and related radar studies use game-based learning for frequency-agile and waveform decisions. These studies establish adaptive radar--jammer interaction as a useful game formulation, but their principal outcomes are link-level quantities such as SINR, detection probability, or waveform utility. FluxPhased complements this line with a mission-level endpoint: a target is successful only when the appropriate service detects the appropriate bearing within a deadline window \cite{heinrich2016deep}.

\subsection{Multifunction-radar resource management}
A separate literature treats multifunction-radar scheduling as an MDP or constrained MDP. Deep Q-learning, actor--critic methods, PPO, SAC, and model-based planning have been applied to scan--track allocation, congested task scheduling, and time or power budgeting. These approaches address the central scheduling trade-off that a radar cannot attend to every task at every instant. They generally use fixed or scripted interference conditions, however, and therefore do not measure how a learning radar changes when the adversary co-adapts. Our benchmark retains the scheduling structure of this literature while adding a learning jammer team and an explicit natural-floor control.

\subsection{Self-play, specialization, and leagues}
Self-play can produce strong coordination but can also specialize agents to the distribution of partners or opponents encountered during training \cite{heinrich2016deep,vinyals2019alphastar,wang2020roma}. Role-oriented multi-agent learning relates team performance to emergent division of labor, and population-based training addresses opponent diversity rather than a single opponent trajectory. The hierarchical macro-strategy model for MOBA games likewise separates high-level target allocation from low-level execution and uses cross-agent coordination \cite{wu2018moba}. We borrow these ideas as motivation; R5 tests a minimal opponent-class mixture, not a full population-based league.

\subsection{Position of this work}
FluxPhased lies at the intersection of these communities. Its agents learn link-facing actions, but the reported endpoint is task completion under service, bearing, deadline, and energy constraints. This enables a controlled question that link-level or single-sided scheduling studies do not answer: whether an apparent defensive advantage survives attacker-count scaling. The resulting comparison is designed around three controls: the idle-radar floor, a scripted-radar jammer view, and a co-located-jammer ablation. The controls let us separate absolute loss, adversarial marginal loss, equilibrium stability, and robustness to a changed opponent class.
